This appendix collects compact reference material used throughout the book. It is organized as a set of self-contained guides and tables intended for quick consultation rather than extended exposition.

\begin{itemize}
  \item \textbf{Appendix A: Post's lattice cheat sheet and completeness tables.} A reference summary of the principal clone-theoretic classes, their closure properties, and the standard completeness criteria used in the Boolean setting.
  \item \textbf{Appendix B: Operad and PROP identities, with coherence diagrams.} A compact list of the identities, interchange laws, and diagrammatic coherence principles that recur in the operadic and PROP-based constructions.
  \item \textbf{Appendix C: ZX, traced monoidal, and decorated cospan guides.} A practical guide to the notation, rewrite rules, and categorical conventions used for ZX-calculus, traced monoidal structure, and decorated cospan formalisms.
  \item \textbf{Appendix D: Metrics catalog for $\varepsilon(r)$.} A catalog of metric definitions, units, and calibration conventions for the error and resolution measures denoted by $\varepsilon(r)$, including the bookkeeping needed to compare values across chapters and build artifacts.
\end{itemize}

The appendices are meant to support cross-referencing and reproducibility. When a later chapter invokes a lattice class, an operadic law, a traced feedback identity, or a metric convention, the corresponding reference material should be located here before consulting the main text.

For consistency with the rest of the book, notation in these appendices follows the conventions established in the front matter. Where a result is stated in abbreviated form, the surrounding chapters provide the full derivation or proof sketch; the appendix entry serves as the canonical lookup location for the statement itself.

\paragraph{Usage note.} These materials are intentionally terse. They are designed to be stable reference points for proofs-as-builds, implementation checks, and calibration workflows, and they may be expanded in future revisions as the surrounding chapters evolve.

\paragraph{Appendix A: Post's lattice cheat sheet and completeness tables.} This appendix records the standard Boolean clone classes, their defining closure properties, and the usual completeness criteria used to recognize functionally complete sets. It is intended as a lookup table for the lattice-theoretic terminology used in the chapters on closure and clones.

\paragraph{Appendix B: Operad and PROP identities.} This appendix summarizes the basic operadic identities, the PROP interchange law, and the coherence patterns that justify diagrammatic manipulations. It is the reference point for equations and string-diagram equalities that are used without rederivation in the main text.

\paragraph{Appendix C: ZX, traced monoidal, and decorated cospan guides.} This appendix collects the notation and rewrite conventions for ZX-calculus, the axioms and trace identities for traced monoidal categories, and the bookkeeping conventions for decorated cospans. It is meant to make the categorical and diagrammatic chapters easier to read and verify.

\paragraph{Appendix D: Metrics catalog for $\varepsilon(r)$.} This appendix lists the metric definitions, units, normalization choices, and calibration procedures associated with $\varepsilon(r)$. It provides the comparison conventions needed to interpret error values consistently across chapters, experiments, and build artifacts.

\begin{itemize}
  \item The appendices are reference-oriented rather than tutorial-oriented.
  \item Each entry is intended to be stable under later chapter revisions.
  \item Cross-references should point here whenever a definition, identity, or calibration rule is used repeatedly.
\end{itemize}

\paragraph{Appendix A: Post's lattice cheat sheet; completeness tables.} The Boolean reference tables in this appendix summarize the standard clone classes appearing in Post's lattice, together with the closure properties that characterize them. The tables are intended to support quick checks of whether a given set of Boolean operations is closed under composition, preserves a relation, or satisfies one of the standard completeness criteria.

\paragraph{Appendix B: Operad/PROP identities; coherence diagrams.} This appendix records the identities and coherence principles that are used repeatedly in the operad and PROP chapters. In particular, it serves as a compact reference for interchange laws, associativity and unit constraints, and the diagrammatic equalities that justify rewriting composite operations.

\paragraph{Appendix C: ZX, traced monoidal, decorated cospan guides.} This appendix gathers the working conventions for ZX-calculus, traced monoidal categories, and decorated cospans. It includes the notation and rewrite rules needed to interpret string diagrams, the feedback identities used in traced settings, and the compositional bookkeeping used for decorated cospan constructions.

\paragraph{Appendix D: Metrics catalog for $\varepsilon(r)$: defs, units, calibration.} This appendix catalogs the metric definitions associated with $\varepsilon(r)$, together with their units, normalization choices, and calibration procedures. The intent is to make comparisons across chapters and build artifacts unambiguous by fixing the conventions used to report and interpret error values.

\begin{itemize}
  \item Definitions are stated in a lookup-friendly form.
  \item Units and normalization choices are recorded alongside each metric.
  \item Calibration conventions are included so that reported values can be compared consistently.
\end{itemize}

\paragraph{Reference scope.} Appendix A supports the Boolean closure and completeness chapters; Appendix B supports the operad, PROP, and diagrammatic coherence chapters; Appendix C supports the ZX, traced, and cospan chapters; and Appendix D supports the metric and calibration conventions used throughout the book. Together, these appendices provide the canonical lookup layer for notation, identities, and measurement practice.
